\documentclass[12pt,a4paper]{book}
\usepackage[utf8]{inputenc}
\usepackage[indonesian]{babel}
\usepackage{geometry}
\geometry{margin=1in}
\usepackage{tocloft}
\usepackage{hyperref}
\usepackage{amsmath}
\usepackage{amssymb}

\title{Prediksi Perubahan Rezim Pasar dengan Machine Learning:\\
Pendekatan Komputasional untuk Mendeteksi Bullish dan Bearish Market}
\author{[Nama Penulis]}
\date{\today}

\begin{document}

\maketitle

\tableofcontents

\chapter*{Prakata}
\addcontentsline{toc}{chapter}{Prakata}
Pasar keuangan adalah binatang yang terus berubah. Apa yang berhasil hari ini mungkin akan gagal total besok. Bagi para pedagang dan investor, kemampuan untuk mengenali "musim" pasar---apakah kita sedang berada di tengah badai penurunan, hari cerah kenaikan, atau kabut ketidakpastian yang mendatar---adalah keterampilan yang membedakan antara mereka yang bertahan dan mereka yang tergerus oleh volatilitas.

Buku ini ditulis dari frustrasi dan harapan. Frustrasi melihat banyaknya model trading yang kaku dan gagal beradaptasi saat kondisi pasar berubah. Harapan bahwa dengan kemajuan teknologi, khususnya \textit{Machine Learning}, kita sekarang memiliki alat untuk mendeteksi perubahan rezim ini secara lebih objektif dan cepat daripada sebelumnya.

Tujuan saya menulis buku ini bukan untuk memberikan "kotak hitam" ajaib penghasil uang, melainkan untuk memberikan kerangka kerja komputasional. Bagaimana kita bisa mengajarkan komputer untuk mengenali pola-pola halus yang menandakan pergeseran sentimen pasar? Bagaimana kita memvalidasi sinyal tersebut? Dan yang terpenting, bagaimana kita menyesuaikan strategi kita berdasarkan deteksi tersebut?

Semoga buku ini menjadi panduan bermanfaat bagi siapa saja yang ingin menavigasi kompleksitas pasar modern dengan bantuan kecerdasan buatan.

\vspace{1cm}
\noindent
[Kota, Tanggal]\\
Penulis

\chapter*{Pengantar}
\addcontentsline{toc}{chapter}{Pengantar}
\section*{Latar Belakang}
Model keuangan tradisional seringkali berasumsi bahwa parameter pasar (seperti rata-rata imbal hasil dan volatilitas) adalah konstan sepanjang waktu atau mengikuti distribusi normal yang dapat diprediksi. Namun, realitas sejarah menunjukkan sebaliknya: pasar sering mengalami perubahan struktural mendadak atau "perubahan rezim". Krisis finansial 2008 dan keruntuhan pasar akibat COVID-19 pada 2020 adalah pengingat brutal bahwa strategi yang didasarkan pada data masa damai bisa hancur dalam hitungan hari saat rezim berubah.

\section*{Tujuan Buku}
Buku ini bertujuan untuk:
\begin{itemize}
    \item Menjelaskan konsep rezim pasar secara mendalam, melampaui definisi biner Bull vs Bear sederhana.
    \item Memperkenalkan aplikasi algoritma Machine Learning (dari HMM hingga Deep Learning) khusus untuk tugas deteksi rezim.
    \item Memberikan panduan praktis tentang data preprocessing dan feature engineering yang relevan untuk data finansial yang \textit{noisy}.
    \item Menjembatani teori keuangan modern dengan praktik data science terapan.
\end{itemize}

\section*{Struktur Buku}
Buku ini dibagi menjadi lima bagian utama:
\begin{itemize}
    \item \textbf{Bagian I: Fondasi Pasar dan Rezim} membahas teori dasar pasar dan mengapa rezim terbentuk.
    \item \textbf{Bagian II: Data dan Feature Engineering} fokus pada seni mengolah data pasar menjadi sinyal yang berguna.
    \item \textbf{Bagian III: Metodologi dan Algoritma} mengupas berbagai model ML untuk klasifikasi rezim.
    \item \textbf{Bagian IV: Evaluasi dan Optimisasi} membahas cara menguji dan memperbaiki model trading.
    \item \textbf{Bagian V: Implementasi dan Aplikasi} memberikan contoh kode dan strategi konkret.
\end{itemize}

\section*{Kepada Siapa Buku Ini Ditujukan}
\begin{itemize}
    \item \textbf{Quantitative Traders/Researchers}: Yang ingin meningkatkan model mereka dengan overlay deteksi rezim.
    \item \textbf{Data Scientists}: Yang tertarik terjun ke domain finansial dan memahami tantangan unik data pasar.
    \item \textbf{Portfolio Managers}: Yang membutuhkan pendekatan sistematis untuk manajemen risiko dinamis.
    \item \textbf{Mahasiswa Keuangan/Komputer}: Yang mencari referensi interdisipliner tentang Financial Machine Learning.
\end{itemize}

\part{Fondasi Pasar dan Rezim}

\chapter{Memahami Rezim Pasar}
\section{Konsep Rezim Pasar}
\subsection{Definisi Rezim Pasar}
Dalam dinamika sistem kompleks seperti pasar keuangan, "rezim" merujuk pada periode stabilitas relatif dalam perilaku statistik data. Jika kita memandang pasar sebagai sistem stokastik, rezim pasar adalah keadaan (\textit{state}) di mana parameter distribusi probabilitas yang mendasari (seperti mean $\mu$, varians $\sigma^2$, dan korelasi $\rho$) cenderung konstan atau berubah secara perlahan.

Berbeda dengan pandangan linier yang menganggap pasar bergerak dalam satu kontinum yang mulus, konsep rezim mengakui adanya perubahan struktural (\textit{structural breaks}). Misalnya, perilaku harga saham saat krisis likuiditas global (Rezim Krisis) secara fundamental berbeda aturan mainnya dibandingkan saat ekonomi tumbuh stabil (Rezim Pertumbuhan). Korelasi antar aset yang biasanya rendah bisa tiba-tiba melonjak mendekati 1 saat kepanikan melanda, menggagalkan strategi diversifikasi tradisional. Memahami pergeseran ini adalah kunci untuk bertahan hidup di pasar.

\subsection{Bull Market vs Bear Market}
Secara tradisional, pelaku pasar membagi kondisi menjadi dua kategori biner:
\begin{itemize}
    \item \textbf{Bull Market}: Didefinisikan secara umum sebagai periode kenaikan harga aset sebesar 20\% atau lebih dari titik terendah sebelumnya. Karakteristik psikologisnya adalah optimisme, kepercayaan investor yang tinggi, dan ekspektasi pertumbuhan ekonomi yang kuat. Secara historis, Bull Market cenderung berlangsung lebih lama (rata-rata 2-5 tahun) dengan volatilitas yang relatif rendah. Data historis S\&P 500 menunjukkan bahwa Bull Market membangun kekayaan secara perlahan namun konsisten.
    \item \textbf{Bear Market}: Didefinisikan sebagai penurunan harga aset sebesar 20\% atau lebih dari puncaknya. Ini ditandai dengan pesimisme luas, ketakutan, dan kontraksi ekonomi. Bear Market cenderung lebih pendek (rata-rata 8-18 bulan) tetapi jauh lebih ganas dan volatil. Penurunan harga sering terjadi sangat cepat ("Pasar naik lewat tangga, turun lewat lift"), yang mencerminkan panik penjualan.
\end{itemize}

Namun, definisi persentase statis ini memiliki kelemahan fatal: ia bersifat \textit{lagging}. Pada saat pasar dikonfirmasi turun 20\%, investor mungkin sudah kehilangan sebagian besar modalnya. Pendekatan Machine Learning mencoba mendeteksi probabilitas peralihan ke Bear Market jauh sebelum ambang batas 20\% itu tercapai.

\subsection{Sideways/Range-bound Market}
Seringkali diabaikan dalam dikotomi Bull/Bear, rezim \textit{Sideways} adalah kondisi di mana harga berosilasi dalam rentang horizontal yang terbatas tanpa tren arah yang jelas. Ini adalah "kuburan" bagi strategi \textit{trend-following} yang mengandalkan momentum berkelanjutan. Sebaliknya, strategi \textit{mean-reversion} (beli di \textit{support}, jual di \textit{resistance}) justru memanen profit terbesar di sini. Mengidentifikasi rezim ini sangat krusial untuk menonaktifkan algoritma yang rentan terhadap sinyal palsu (\textit{whipsaw}).

\subsection{Volatility Regimes}
Selain arah (Directional Regime), dimensi kedua yang tak kalah penting adalah volatilitas. Kita bisa membagi pasar menjadi:
\begin{itemize}
    \item \textbf{Low Volatility Regime}: Pergerakan harga harian kecil dan teratur. Sering terjadi pada fase pemulihan ekonomi atau Bull Market yang matang. Strategi dengan leverage tinggi (\textit{carry trade}) sering bekerja baik di sini.
    \item \textbf{High Volatility Regime}: Pergerakan harga harian besar dan tak terduga. Ini biasanya terjadi saat rilis berita makroekonomi besar, krisis geopolitik, atau transisi menuju Bear Market. Risiko \textit{gap opening} sangat tinggi, menuntut manajemen risiko yang ketat.
\end{itemize}

\subsection{Market Cycles dan Fase}
Richard Wyckoff, seorang pionir analisis teknikal awal abad ke-20, memetakan siklus pasar menjadi empat fase distrik yang relevan hingga hari ini:
\begin{enumerate}
    \item \textbf{Akumulasi}: Periode sideways setelah penurunan panjang, di mana \textit{smart money} (institusi) mulai membeli aset secara diam-diam dari investor ritel yang putus asa.
    \item \textbf{Markup (Uptrend)}: Fase kenaikan harga yang nyata, partisipasi publik mulai masuk, dan optimisme meningkat. Ini adalah fase Bull Market.
    \item \textbf{Distribusi}: Periode sideways di puncak, di mana \textit{smart money} mulai menjual kepemilikan mereka kepada investor ritel yang baru masuk karena eforia (FOMO).
    \item \textbf{Markdown (Downtrend)}: Fase penurunan harga saat pasokan melebihi permintaan, memicu Bear Market.
\end{enumerate}
Machine Learning dapat dilatih untuk mengklasifikasikan data pasar saat ini ke dalam salah satu dari empat fase Wyckoff ini, memberikan peta jalan taktis bagi trader.

\section{Pentingnya Identifikasi Rezim}
\subsection{Asset Allocation Strategy yang Dinamis}
Model portofolio statis seperti portofolio 60/40 (60\% saham, 40\% obligasi) berasumsi bahwa korelasi negatif antara saham dan obligasi akan selalu melindungi investor. Namun, pada tahun 2022, baik saham maupun obligasi jatuh bersamaan karena rezim inflasi tinggi.
Dengan identifikasi rezim yang akurat, manajer investasi dapat melakukan \textit{Dynamic Asset Allocation}:
\begin{itemize}
    \item Saat deteksi \textbf{Rezim Inflasi}: Alihkan dana ke Komoditas atau Emas.
    \item Saat deteksi \textbf{Rezim Resesi}: Pindah ke *Cash* atau *Government Bonds* jangka pendek.
    \item Saat deteksi \textbf{Rezim Pertumbuhan}: *Overweight* pada Ekuitas (Saham).
\end{itemize}

\subsection{Risk Management Adaptif}
Manajemen risiko yang kaku seringkali gagal. Stop-loss 2\% mungkin terlalu sempit di pasar volatil (mudah terkena \textit{stop-out} prematur) tetapi terlalu lebar di pasar tenang.
Sistem berbasis rezim menyesuaikan parameter risiko secara otomatis:
\begin{itemize}
    \item \textbf{Position Sizing}: Kurangi ukuran posisi setengahnya saat model mendeteksi probabilitas tinggi beralih ke rezim volatilitas tinggi.
    \item \textbf{Volatility-Adjusted Stop Loss}: Perlebar jarak stop-loss saat volatilitas meningkat untuk memberi ruang bernafas bagi harga (*atr-based stops*).
\end{itemize}

\subsection{Trading Strategy Selection}
Tidak ada strategi tunggal yang berhasil di semua kondisi pasar.
\begin{itemize}
    \item \textbf{Trend Follower}: Profit besar di Bull/Bear Market, rugi besar di Sideways.
    \item \textbf{Mean Reversion}: Profit konsisten di Sideways, hancur lebur saat pasar trending kuat (melawan tren).
\end{itemize}
Sistem meta-learning dapat bertindak sebagai "pelatih" yang memilih pemain (strategi) mana yang diterjunkan ke lapangan berdasarkan kondisi cuaca (rezim) saat ini.

\section{Sejarah Market Regimes: Studi Kasus}
Mempelajari sejarah bukan sekadar hafalan, tetapi memahami pola berulang yang didorong oleh psikologi manusia yang tetap sama.

\subsection{Great Depression (1929-1939)}
Contoh paling ekstrem dari Bear Market sekuler. Pasar saham AS kehilangan hampir 90\% nilainya. Ini mengajarkan bahwa pasar tidak selalu "kembali normal" dalam waktu singkat. Rezim deflasi yang parah mengubah korelasi aset secara drastis, di mana *Cash is King*. Perubahan regulasi pasca-depresi menciptakan rezim pasar baru yang bertahan puluhan tahun.

\subsection{Dot-com Bubble (1995-2002)}
Studi kasus klasik tentang eforia irasional. Valuasi saham teknologi terputus dari fundamental (banyak perusahaan tanpa profit dihargai miliaran dolar). Rezim ini ditandai oleh volatilitas sisi atas (\textit{upside volatility}) yang ekstrem. Pelajaran utamanya: pasar bisa tetap irasional lebih lama daripada kemampuan Anda menahan solvabilitas (seperti kata John Maynard Keynes). Deteksi rezim berbasis valuasi akan memberi sinyal "Bubble" bertahun-tahun sebelum pecah.

\subsection{Financial Crisis 2008}
Krisis ini unik karena sifatnya yang sistemik. Produk derivatif kompleks (MBS/CDO) menyembunyikan risiko yang sebenarnya. Saat gelembung perumahan pecah, efek penularan (\textit{contagion}) menyebar ke seluruh dunia. Rezim ini ditandai oleh krisis likuiditas: aset bagus pun dijual paksa untuk menutupi *margin call*. Korelasi antar semua kelas aset mendekati 1, menggagalkan diversifikasi. Machine Learning yang memantau *Credit Spreads* (selisih yield obligasi korporasi dan pemerintah) bisa mendeteksi tanda-tanda stres sistemik ini berbulan-bulan sebelumnya.

\subsection{COVID-19 Market Crash (2020)}
Contoh \textit{Black Swan} sempurna. Penurunan tercepat dalam sejarah pasar (Bear Market dalam hitungan minggu), diikuti oleh pemulihan berbentuk V tercepat dalam sejarah berkat intervensi bank sentral. Ini menciptakan siklus pasar yang "dipercepat". Model yang dilatih pada data historis lambat mungkin gagal beradaptasi dengan kecepatan perubahan ini, menekankan perlunya model adaptif (\textit{online learning}).

\chapter{Teori Pasar Keuangan}
Landasan teoritis sangat penting untuk membumikan model Machine Learning kita. Tanpa pemahaman tentang ekonomi dan perilaku pasar, seorang data scientist hanya akan menemukan korelasi palsu (\textit{spurious correlations}) dalam data. Kita perlu memahami "mengapa" pasar bergerak (teori) sebelum memprediksi "bagaimana" ia bergerak (ML).

\section{Efficient Market Hypothesis (EMH)}
Diajukan oleh Eugene Fama pada tahun 1960-an, EMH adalah pilar keuangan modern yang menyatakan bahwa harga aset selalu mencerminkan secara penuh semua informasi yang tersedia. Menurut teori ini, tidak mungkin untuk mengalahkan pasar secara konsisten (mencetak \textit{alpha}) karena setiap informasi baru langsung diserap ke dalam harga seketika.

\subsection{Tiga Bentuk Efisiensi}
EMH dibagi menjadi tiga tingkatan:
\begin{enumerate}
    \item \textbf{Weak Form (Bentuk Lemah)}: Harga saat ini sudah mencerminkan semua data harga dan volume masa lalu. Implikasinya: \textit{Analisis Teknikal} (melihat grafik) tidak berguna karena pola masa lalu tidak memprediksi masa depan. Pergerakan harga adalah "Jalan Acak" (\textit{Random Walk}).
    \item \textbf{Semi-strong Form (Bentuk Semi-kuat)}: Harga mencerminkan semua informasi publik (laporan keuangan, berita, data makro). Implikasinya: \textit{Analisis Fundamental} tidak memberikan keunggulan, karena begitu berita bagus rilis, harga sudah menyesuaikan diri sebelum Anda sempat klik "Beli".
    \item \textbf{Strong Form (Bentuk Kuat)}: Harga mencerminkan \textit{semua} informasi, termasuk informasi privat atau orang dalam (\textit{insider}). Implikasinya: Bahkan CEO perusahaan atau insider trading pun tidak bisa mengalahkan pasar dalam jangka panjang.
\end{enumerate}

\subsection{Kritik dan Adaptive Market Hypothesis}
Bagi praktisi ML, EMH dalam bentuk murninya adalah berita buruk---jika pasar efisien sempurna, maka tidak ada pola untuk dipelajari, dan prediksi adalah mustahil. Namun, segunung bukti empiris dan kesuksesan dana lindung nilai (\textit{hedge funds}) kuantitatif membantah EMH bentuk kuat.
Andrew Lo mengusulkan **Adaptive Market Hypothesis (AMH)** sebagai sintesis. AMH menyatakan bahwa efisiensi pasar bukanlah kondisi biner (Efisien/Tidak Efisien), melainkan spektrum yang berubah-ubah tergantung kondisi lingkungan (rezim).
\begin{itemize}
    \item Di pasar yang tenang dan matang, efisiensi tinggi (sulit diprediksi).
    \item Di masa panik, krisis, atau adanya pelaku pasar baru yang tidak rasional, efisiensi menurun, menciptakan inefisiensi yang bisa dieksploitasi oleh algoritma ML.
\end{itemize}
Buku ini berlandaskan pada premis AMH: kita menggunakan ML untuk mencari "kantong-kantong inefisiensi" ini.

\section{Behavioral Finance}
Jika EMH mengasumsikan pelaku pasar adalah robot rasional ("Homo Economicus"), Behavioral Finance mengakui bahwa pasar diisi oleh manusia yang emosional dan bias. Perubahan rezim seringkali lebih didorong oleh psikologi massa daripada perubahan fundamental ekonomi.

\subsection{Market Psychology: Fear and Greed}
Pasar berayun seperti pendulum antara ketakutan (fear) dan keserakahan (greed).
\begin{itemize}
    \item \textbf{Greed (Rezim Bull)}: Investor mengabaikan risiko, valuasi menjadi tidak masuk akal, dan rasa takut ketinggalan (FOMO) mendorong harga naik melampaui nilai wajarnya.
    \item \textbf{Fear (Rezim Bear)}: Investor mengabaikan peluang, menjual aset bagus dengan harga diskon karena panik. VIX (indeks volatilitas) adalah "termometer ketakutan" yang bisa menjadi fitur input ML yang sangat kuat.
\end{itemize}

\subsection{Bias Kognitif dalam Trading}
\begin{itemize}
    \item \textbf{Herding Behavior (Perilaku Menggembala)}: Manusia merasa lebih aman mengikuti orang banyak. Jika semua orang menjual, insting kita menyuruh ikutan jual. Ini menciptakan tren momentum yang panjang dan berlebihan (\textit{overshoot}).
    \item \textbf{Loss Aversion}: Teori Prospek (Kahneman & Tversky) menunjukkan bahwa rasa sakit akibat rugi \$100 dua kali lebih besar daripada rasa senang untung \$100. Ini menjelaskan mengapa pasar jatuh lebih cepat daripada naiknya (Panic Selling lebih intens daripada Panic Buying).
    \item \textbf{Confirmation Bias}: Trader cenderung hanya mencari berita yang mendukung posisi mereka dan mengabaikan sinyal bahaya. Algoritma ML yang objektif membantu memitigasi bias manusia ini.
\end{itemize}

\section{Market Microstructure}
Sementara teori makro menjelaskan tren jangka panjang, *Market Microstructure* membahas mekanika detail bagaimana perdagangan terjadi di level *tick-by-tick*. Untuk strategi frekuensi tinggi atau eksekusi algoritma, pemahaman ini vital.

\subsection{Likuiditas dan Order Book}
Pasar terdiri dari \textit{Limit Order Book} (LOB).
\begin{itemize}
    \item \textbf{Bid}: Harga tertinggi yang bersedia dibayar pembeli.
    \item \textbf{Ask/Offer}: Harga terendah yang diminta penjual.
    \textbf{Spread}: Selisih Ask - Bid. Spread yang melebar adalah indikator awal stres pasar atau penurunan likuiditas.
\end{itemize}
Dalam rezim krisis, likuiditas sering "menguap" tiba-tiba, memperburuk kejatuhan harga karena tidak ada pembeli yang bersedia menampung ( *liquidity black hole* ).

\subsection{Peran High-Frequency Trading (HFT)}
Sebagian besar volume pasar modern didominasi oleh algoritma HFT. Mereka menyediakan likuiditas tetapi juga bisa menariknya dalam milidetik (*quote stuffing/spoofing*). Pola-pola mikro yang dihasilkan HFT menciptakan "noise" yang harus dibersihkan sebelum melatih model ML untuk prediksi trend jangka panjang. Namun, jejak kaki HFT (seperti ketidakseimbangan order flow) juga bisa menjadi fitur prediktif jangka pendek.

\section{Analisis Teknikal sebagai Fitur ML}
Banyak akademisi mencemooh analisis teknikal sebagai "voodoo", tetapi bagi praktisi ML, indikator teknikal adalah metode \textit{feature engineering} yang valid untuk menangkap psikologi pasar.
\begin{itemize}
    \item \textbf{Support/Resistance}: Level memori pasar di mana harga pernah berbalik.
    \item \textbf{Moving Averages}: Filter low-pass yang meredam noise untuk menonjolkan tren sinyal.
    \item \textbf{RSI/Oscillators}: Ukuran kecepatan perubahan harga untuk mendeteksi kondisi jenuh (overbought/oversold).
\end{itemize}
Alih-alih menggunakan aturan kaku ("Beli jika RSI < 30"), ML menggunakan nilai-nilai ini sebagai input vektor untuk mempelajari aturan keputusan yang lebih kompleks dan nuansa secara non-linear.

\chapter{Pengantar Machine Learning untuk Finance}
Machine Learning (ML) telah mengubah lanskap industri keuangan, bergerak dari sekadar alat statistik tambahan menjadi mesin utama penggerak keputusan investasi. Namun, menerapkan ML pada data keuangan bukanlah tugas yang mudah; bidang ini penuh dengan jebakan yang tidak ditemukan dalam domain lain seperti pengenalan gambar atau pemrosesan bahasa alami.

\section{Mengapa Machine Learning untuk Pasar?}
\subsection{Kompleksitas Data Finansial}
Data pasar keuangan adalah representasi dari psikologi jutaan pelaku pasar yang berinteraksi secara real-time. Data ini bersifat \textit{noisy}, \textit{chaotic}, dan seringkali tidak terstruktur. Model ekonometrika tradisional (seperti regresi linier OLS) bekerja berdasarkan asumsi yang seringkali tidak realistis, seperti hubungan yang linier dan residu yang terdistribusi normal.
Machine Learning, khususnya model non-parametrik seperti \textit{Decision Trees} dan \textit{Neural Networks}, tidak memaksakan asumsi bentuk pada data. Mereka memiliki fleksibilitas untuk mempelajari pola non-linear yang kompleks dan interaksi antar variabel yang halus (misalnya: korelasi antara Emas dan Dolar mungkin hanya kuat saat inflasi di atas 3%, dan ML bisa menangkap nuansa ini secara otomatis).

\subsection{Alternative Data Integration}
Revolusi terbesar ML di finance adalah kemampuannya menelan data non-tradisional. Sementara model lama terbatas pada harga dan volume, ML dapat memproses:
\begin{itemize}
    \textbf{Text Data}: Sentimen berita, tweet, dan transkrip earning calls.
    \item \textbf{Image Data}: Citra satelit parkiran Walmart untuk memprediksi pendapatan ritel.
    \item \textbf{Network Data}: Hubungan rantai pasok antar perusahaan global.
\end{itemize}
Kemampuan untuk mengubah data mentah yang tidak terstruktur ini menjadi sinyal trading yang dapat ditindaklanjuti adalah keunggulan kompetitif utama (\textit{edge}).

\section{Paradigma Machine Learning dalam Konteks Trading}
\subsection{Supervised Learning: Prediksi Arah}
Ini adalah paradigma yang paling umum digunakan. Tujuannya adalah memetakan fitur (X) ke target (Y).
\begin{itemize}
    \item \textbf{Klasifikasi}: Memprediksi arah harga diskrit. Contoh: "Apakah harga besok akan Naik (1) atau Turun (0)?". Aplikasi lain adalah deteksi penipuan kartu kredit.
    \item \textbf{Regresi}: Memprediksi nilai kontinu. Contoh: "Berapa harga penutupan S\&P 500 besok?". Perlu dicatat bahwa di finance, memprediksi \textit{return} (perubahan harga) jauh lebih stabil daripada memprediksi harga absolut.
\end{itemize}

\subsection{Unsupervised Learning: Penemuan Pola}
Di sini, kita tidak memiliki label target (Y). Algoritma dibiarkan mencari struktur tersembunyi dalam data.
\begin{itemize}
    \item \textbf{Clustering}: Mengelompokkan saham yang bergerak serupa untuk diversifikasi portofolio. Atau, seperti topik utama buku ini: mengelompokkan hari-hari perdagangan ke dalam "Rezim Pasar" (misal: Cluster Volatilitas Tinggi vs Rendah).
    \item \textbf{Dimensionality Reduction (PCA)}: Mengompres ribuan indikator teknikal menjadi beberapa "Faktor Risiko" utama untuk menghindari kutukan dimensi.
\end{itemize}

\subsection{Reinforcement Learning: Optimalisasi Tindakan}
Paradigma ini paling mirip dengan cara manusia belajar trading. Sebuah "Agen" berinteraksi dengan "Lingkungan" (Pasar) dan menerima "Reward" (Profit/Loss). Tujuannya bukan sekadar memprediksi harga, tetapi mengambil serangkaian keputusan (Beli/Jual/Tahan) yang memaksimalkan total kekayaan jangka panjang. RL sangat kuat untuk manajemen risiko dan eksekusi order yang optimal.

\section{Tantangan Unik ML di Finance}
Mengapa fisikawan pintar sering gagal menjadi trader sukses? Karena data keuangan berbeda.

\subsection{Low Signal-to-Noise Ratio}
Dalam pengenalan gambar, gambar kucing tetaplah kucing; sinyalnya jelas. Di pasar, sinyal pergerakan harga seringkali sangat lemah, tersembunyi di balik gunungan \textit{noise} (transaksi acak). Akurasi prediksi 51-54\% sudah dianggap "Sangat Bagus" di trading frekuensi tinggi. Ini menuntut disiplin statistik yang ketat untuk tidak menipu diri sendiri.

\subsection{Non-stationarity dan Concept Drift}
Ini adalah musuh terbesar. Distribusi data berubah seiring waktu. Strategi yang dilatih pada tahun 2010 (pasar tenang pasca-krisis) mungkin gagal total di 2020 (krisis pandemi). Model ML di finance harus bersifat \textit{adaptif} atau sering dilatih ulang (\textit{retrained}) untuk tetap relevan. Apa yang kita pelajari hari ini mungkin kedaluwarsa besok.

\subsection{Backtest Overfitting}
"Menyiksa data sampai dia mengaku." Sangat mudah untuk membuat strategi yang menghasilkan Sharpe Ratio 3.0 di data masa lalu (\textit{in-sample}) dengan memilih parameter terbaik secara surut. Namun, strategi ini hampir pasti rugi di pasar nyata (\textit{out-of-sample}). Teknik validasi khusus diperlukan untuk memitigasi ini.

\section{Validasi Model yang Benar}
\subsection{Bahaya Acak (Random K-Fold)}
Dalam ML standar, kita sering mengacak data sebelum membaginya menjadi K-Fold. Di data time series, ini fatal! Menggunakan data masa depan (esok) untuk memprediksi masa lalu (kemarin) adalah kecurangan (\textit{look-ahead bias}).

\subsection{Walk-Forward Analysis}
Metode validasi emas untuk time series.
\begin{enumerate}
    \item Latih model pada data Tahun 1-3. Uji di Tahun 4.
    \item Geser jendela: Latih di Tahun 2-4. Uji di Tahun 5.
    \item Dan seterusnya.
\end{enumerate}
Ini mensimulasikan pengalaman trader yang sebenarnya yang harus beradaptasi dengan aliran data baru.

\subsection{Purged K-Fold Cross Validation}
Diusulkan oleh Marcos Lopez de Prado. Teknik ini menghapus data di perbatasan antara set Training dan Testing (purging) untuk mencegah kebocoran informasi dari korelasi serial (karena trade hari ini dipengaruhi trade kemarin).

\part{Data dan Feature Engineering}

\chapter{Sumber Data Pasar}
Bahan bakar utama algoritma ML adalah data. Kualitas output model (Garbage Out) sangat bergantung pada kualitas input data (Garbage In).

\section{Price Data (Market Data)}
\subsection{OHLCV: Standar Industri}
Format data yang paling umum adalah \textbf{Open, High, Low, Close, Volume (OHLCV)} yang di-sampel berdasarkan waktu (Time Bars).
\begin{itemize}
    \item \textbf{Open}: Harga transaksi pertama pada periode tersebut.
    \item \textbf{High/Low}: Rentang ekstrem harga, penting untuk mengukur volatilitas.
    \item \textbf{Close}: Harga terakhir, sering dianggap sebagai harga "konsensus" final periode itu.
    \item \textbf{Volume}: Jumlah aset yang berpindah tangan, indikator validitas pergerakan harga.
\end{itemize}
Kelemahan Time Bars adalah mereka memotong aktivitas pasar secara artifisial. Satu jam di jam makan siang yang sepi diperlakukan sama dengan satu jam saat rilis berita Non-Farm Payroll, padahal nilai informasinya sangat berbeda.

\subsection{Tick Data}
Sumber data bervariasi dari yang gratis (Yahoo Finance) hingga premium (Bloomberg, Refinitiv). Kualitas data berbayar biasanya jauh lebih baik dalam hal penyesuaian aksi korporasi dan kelengkapan historis.

\section{Fundamental Data}
\subsection{Economic Indicators}
Data makroekonomi (NFP, GDP, CPI) yang dirilis secara berkala. Perubahan rezim sering dipicu oleh kejutan pada rilis data ini.

\subsection{Earnings Reports}
Laporan laba rugi kuartalan perusahaan. Kejutan pendapatan (\textit{earnings surprise}) adalah penggerak harga utama untuk saham individu.

\subsection{Balance Sheet Metrics}
Rasio utang, arus kas, dan nilai buku. Penting untuk model investasi jangka panjang tipe \textit{value investing}.

\subsection{Macroeconomic Data}
Suku bunga bank sentral dan kebijakan moneter adalah pendorong utama siklus makro pasar.

\section{Alternative Data}
\subsection{Sentiment dari News dan Social Media}
Analisis sentimen Twitter/X, Reddit (WallStreetBets), dan berita keuangan menggunakan NLP. "Fear" di media sosial sering mendahului penurunan harga.

\subsection{Web Scraping}
Mengumpulkan data harga produk dari situs e-commerce untuk memprediksi inflasi secara real-time, atau data lowongan kerja untuk memprediksi kesehatan ekonomi.

\subsection{Satellite Imagery}
Menghitung jumlah mobil di parkiran retail untuk memprediksi pendapatan, atau memantau tangki minyak untuk mengestimasi cadangan global.

\subsection{Credit Card Transaction}
Data agregat pengeluaran konsumen memberikan pandangan hampir real-time tentang status ekonomi ritel.

\subsection{Google Trends}
Volume pencarian kata kunci tertentu (misal: "recesi", "beli emas", "cara jual saham") berkorelasi kuat dengan perilaku investor ritel.

\section{Market Microstructure Data}
\subsection{Order Book Data (L2/L3)}
Melihat kedalaman pasar (\textit{Market Depth}). Ketidakseimbangan pesanan beli/jual di order book (\textit{Order Imbalance}) adalah prediktor kuat pergerakan harga jangka pendek.

\subsection{Trade Flow}
Arah agresivitas perdagangan. Apakah pembeli memukul harga ask (agresif) atau penjual memukul harga bid?

\subsection{Price Discovery}
Proses bagaimana harga baru terbentuk melalui interaksi pembeli dan penjual.

\subsection{High-Frequency Trading (HFT)}
Algoritma ultra-cepat yang menyediakan likuiditas tetapi juga dapat memicu volatilitas kilat. Model kita perlu memperhitungkan noise yang dihasilkan HFT.

\chapter{Data Preprocessing}
\section{Data Cleaning}
\subsection{Missing Data Handling}
Dalam time series keuangan, menghapus baris (\textit{drop}) dapat merusak struktur waktu. Metode umum adalah \textit{Forward Fill} (menggunakan harga terakhir yang diketahui) atau interpolasi linier untuk gap pendek.

\subsection{Outlier Detection}
\textit{Flash crash} atau kesalahan data feed dapat menciptakan \textit{spikes} harga yang tidak realistis. Filter statistik seperti \textit{Median Absolute Deviation} (MAD) lebih kuat daripada Standar Deviasi untuk mendeteksi ini.

\subsection{Corporate Actions Adjustment}
Harga saham harus disesuaikan (\textit{adjusted}) untuk dividen dan pemecahan saham (\textit{stock splits}) agar grafik sejarah konsisten. Tanpa ini, model akan melihat "gap" harga palsu sebagai sinyal trading.

\subsection{Data Quality Checks}
Memastikan tidak ada duplikasi cap waktu (\textit{timestamps}), harga negatif, atau volume nol pada jam perdagangan aktif.

\section{Normalization dan Transformation}
\subsection{Returns vs Prices}
Harga aset mentah biasanya non-stasioner (memiliki tren). Model ML bekerja lebih baik pada data stasioner. Menggunakan \textit{returns} (persentase perubahan) membuat data lebih stasioner.

\subsection{Log Returns}
Lebih disukai daripada \textit{simple returns} karena sifat aditifnya (log return periode panjang adalah jumlah log return periode pendek) dan distribusi yang lebih simetris (mendekati normal), memudahkan pemodelan statistik.

\subsection{Standardization}
Banyak algoritma (seperti SVM dan Neural Networks) sensitif terhadap skala. \textit{Z-score normalization} (mean=0, std=1) atau \textit{Min-Max Scaling} memastikan semua fitur memiliki bobot awal yang setara.

\subsection{Handling Non-stationarity}
Selain differencing (returns), teknik seperti \textit{Fractional Differentiation} (Lopez de Prado) dapat digunakan untuk membuat seri stasioner sambil mempertahankan memori jangka panjang sebanyak mungkin.

\section{Resampling}
Mengubah frekuensi data.

\subsection{Time-based Resampling}
Metode standar (bar 1 menit, 1 jam, 1 hari). Kelemahannya: pasar tidak aktif secara merata sepanjang waktu (misalnya, jam makan siang sepi).

\subsection{Tick Bars}
Sampling setiap N transaksi. Ini menyesuaikan dengan aktivitas pasar, tetapi ukuran transaksi bisa bervariasi.

\subsection{Volume Bars}
Sampling setiap N saham/kontrak yang diperdagangkan. Lebih baik dalam menangkap periode aktivitas tinggi.

\subsection{Dollar Bars}
Sampling setiap nilai transaksi dolar tertentu (misal: setiap \$1 juta yang diperdagangkan). Ini adalah metode paling stabil secara statistik karena menormalkan fluktuasi harga dan volume.

\section{Train-Test Split untuk Time Series}
\subsection{Temporal Split}
Data harus dipisah secara kronologis (Masa Lalu untuk Train, Masa Depan untuk Test). Pengocokan acak (\textit{random shuffling}) dilarang keras karena merusak urutan waktu.

\subsection{Purging dan Embargo}
Konsep dari Marcos Lopez de Prado.
\begin{itemize}
    \item \textbf{Purging}: Menghapus data di perbatasan antara Train dan Test set untuk mencegah kebocoran informasi dari label yang "tumpang tindih" (overlapping labels).
    \item \textbf{Embargo}: Menambahkan jarak waktu ekstra setelah set Test sebelum menggunakan data selanjutnya untuk Train (pada cross-validation) guna menghilangkan korelasi serial yang tersisa.
\end{itemize}

\subsection{Cross-Validation Strategies}
Menggunakan \textit{Walk-Forward Validation} atau \textit{Blocking Rolling Window} alih-alih K-Fold standar.

\chapter{Technical Indicators sebagai Features}
Indikator teknikal adalah bentuk paling dasar dari \textit{Feature Engineering} dalam trading.

\section{Trend Indicators}
\subsection{Moving Averages (SMA, EMA)}
Meratakan data harga untuk melihat tren dasar. SMA (Simple) memberi bobot sama rata, EMA (Exponential) memberi bobot lebih pada data terbaru. Crossover (Golden Cross/Death Cross) adalah fitur biner sederhana.

\subsection{MACD (Moving Average Convergence Divergence)}
Indikator momentum mengikuti tren yang menunjukkan hubungan antara dua moving averages dari harga. Histogram MACD sering digunakan untuk mendeteksi perubahan momentum tren.

\subsection{ADX (Average Directional Index)}
Mengukur kekuatan tren tanpa mempedulikan arahnya. Nilai ADX tinggi menunjukkan pasar sedang trending kuat (cocok untuk strategi trend-following), ADX rendah menunjukkan pasar sideways (cocok untuk mean-reversion).

\subsection{Parabolic SAR}
Indikator untuk menentukan titik keluar (stop loss) potensial dan pembalikan arah.

\section{Momentum Indicators}
\subsection{RSI (Relative Strength Index)}
Osilator yang mengukur kecepatan dan perubahan pergerakan harga. Level >70 (Overbought) dan <30 (Oversold) adalah fitur input klasik untuk mendeteksi pembalikan mean.

\subsection{Stochastic Oscillator}
Membandingkan harga penutupan tertentu dengan rentang harganya selama periode waktu tertentu. Berguna untuk menentukan titik masuk di pasar sideways.

\subsection{Rate of Change (ROC)}
Persentase perubahan harga antara periode saat ini dan n periode lalu. Mengukur kecepatan tren secara murni.

\subsection{Williams \%R}
Indikator momentum yang bergerak terbalik dibanding Stochastic, mengukur level overbought/oversold.

\section{Volatility Indicators}
\subsection{Bollinger Bands}
Pita volatilitas yang ditempatkan di atas dan di bawah moving average. Penyempitan pita (\textit{Bollinger Squeeze}) sering mendahului ledakan volatilitas besar (perubahan rezim).

\subsection{ATR (Average True Range)}
Mengukur volatilitas pasar dengan menguraikan seluruh rentang harga aset untuk periode tersebut. Sering digunakan untuk menetapkan target stop-loss dinamis (volatility-based stop).

\subsection{Standard Deviation}
Ukuran statistik dispersi harga dari rata-ratanya.

\subsection{Keltner Channels}
Serupa dengan Bollinger Bands tetapi menggunakan ATR untuk menetapkan lebar pita, sehingga lebih responsif terhadap perubahan volatilitas nyata.

\section{Volume Indicators}
\subsection{On-Balance Volume (OBV)}
Menggunakan aliran volume untuk memprediksi perubahan harga saham. Ide dasarnya: volume mendahului harga. Divergensi antara OBV dan harga adalah sinyal kuat.

\subsection{Volume Price Trend}
Menggabungkan harga dan volume.

\subsection{Accumulation/Distribution}
Mencoba mengukur aliran uang masuk dan keluar dari suatu sekuritas.

\subsection{Money Flow Index}
RSI yang ditimbang dengan volume.

\section{Custom Indicators}
\subsection{Kombinasi Indicators}
Membuat "super-indikator" dengan menggabungkan sinyal dari trend dan momentum (misal: "RSI di atas 50 DAN Harga di atas EMA200").

\subsection{Adaptive Indicators}
Indikator yang parameternya menyesuaikan diri secara otomatis dengan siklus dominan pasar saat ini (misal: \textit{Adaptive Moving Average} karya Perry Kaufman).

\subsection{Multi-timeframe Features}
Memasukkan fitur dari kerangka waktu yang lebih tinggi (misal: tren Mingguan) sebagai konteks untuk model yang memprediksi pergerakan Harian.


\chapter{Advanced Feature Engineering}
Jika indikator teknikal adalah "kulit" pasar, maka fitur statistik dan mikrostruktur adalah "organ dalam"nya. Di bab ini, kita akan menggali fitur-fitur canggih yang jarang digunakan oleh trader ritel tetapi menjadi makanan sehari-hari algoritma institusi.

\section{Statistical Features}
Fitur statistik menangkap properti distribusi data yang tidak terlihat oleh mata telanjang pada grafik harga.

\subsection{Rolling Statistics (Mean, Variance)}
Pasar yang efisien seharusnya memiliki return dengan mean nol. Pergeseran mean (\textit{drift}) yang signifikan dalam jendela waktu rolling (misal 20 hari) adalah sinyal tren. Lebih penting lagi adalah \textit{Rolling Variance}. Lonjakan varians seringkali terjadi \textit{sebelum} harga jatuh tertembus. Kita menggunakan jendela waktu eksponensial (EWM) untuk memberi bobot lebih pada data terbaru.

\subsection{Higher Moments (Skewness & Kurtosis)}
Dua momen pertama (Mean, Varians) mengasumsikan distribusi Normal. Pasar tidak Normal.
\begin{itemize}
    \item \textbf{Skewness (Kemencengan)}: Mengukur asimetri distribusi return. Skewness negatif yang besar (ekor kiri panjang) menandakan risiko \textit{crash} yang tinggi ("Naga Hitam").
    \item \textbf{Kurtosis (Keruncingan)}: Mengukur ketebalan ekor (\textit{fat tails}). Kurtosis > 3 (Leptokurtic) berarti probabilitas kejadian ekstrem (sangat untung atau sangat rugi) jauh lebih tinggi daripada prediksi model Gaussian.
\end{itemize}

\subsection{Autocorrelation dan Serial Correlation}
Mengukur seberapa kuat harga hari ini dipengaruhi oleh harga kemarin.
\[ \rho_k = \frac{\sum_{t=k+1}^T (y_t - \bar{y})(y_{t-k} - \bar{y})}{\sum_{t=1}^T (y_t - \bar{y})^2} \]
Pasar yang efisien memiliki autokorelasi mendekati nol (Random Walk). Lonjakan autokorelasi positif menandakan tren (momentum), sedangkan negatif menandakan mean-reversion.

\subsection{Fractal Dimension (Hurst Exponent)}
Mengukur "kekasaran" grafik harga untuk menentukan sifat memori jangka panjang (Long-term Memory). Nilai Hurst Exponent ($H$):
\begin{itemize}
    \item $0.5 < H < 1.0$: Trending (Persistent). Sinyal untuk strategi Trend Following.
    \item $0.0 < H < 0.5$: Mean Reverting (Anti-persistent). Sinyal untuk strategi Reversion.
    \item $H = 0.5$: Random Walk.
\end{itemize}

\section{Microstructure Features}
Mengintip ke dalam mesin pencocokan order (matching engine).

\subsection{Bid-Ask Spread dynamics}
Spread bukan biaya statis. Spread yang melebar secara tiba-tiba adalah proksi ketakutan \textit{market maker}. Saat market maker mendeteksi adanya "Informed Trader", mereka melebarkan spread untuk melindungi diri. Ini sinyal volatilitas.

\subsection{Order Imbalance}
Rasio volume beli agresif vs volume jual agresif.
\[ \rho = \frac{V_{buy} - V_{sell}}{V_{buy} + V_{sell}} \]
Jika order beli mendominasi ($\rho > 0$) tetapi harga tidak naik, ini tanda \textit{absorption} (distribusi tersembunyi).

\subsection{Kyle's Lambda (Amihud Illiquidity)}
Mengukur \textit{market impact cost}: seberapa besar harga bergerak untuk setiap unit volume yang diperdagangkan ($ \Delta P / \Delta V $). Lonjakan Lambda adalah tanda bahaya kerapuhan pasar (*fragility*).

\section{Image-based Encoding (Time-Series to Image)}
Model modern mulai meninggalkan input vektor 1D dan beralih ke representasi visual untuk memanfaatkan kekuatan Computer Vision.
\begin{itemize}
    \item \textbf{Gramian Angular Field (GAF)}: Mengonversi return 1D menjadi matriks citra 2D dengan mempertahankan korelasi temporal melalui koordinat polar. Ada dua varian: \textit{Gramian Angular Summation Field} (GASF) dan \textit{Gramian Angular Difference Field} (GADF).
    \item \textbf{Markov Transition Field (MTF)}: Mengodekan probabilitas transisi antar-kuantil harga menjadi gambar tekstur.
    \item \textbf{Recurrence Plots}: Memvisualisasikan kapan suatu kondisi harga "berulang" di masa depan, sangat efektif untuk mendeteksi rezim siklikal.
\end{itemize}

\section{Feature Selection Techniques}
\subsection{Correlation Clustering}
Mengelompokkan fitur statistik yang redundan dengan Hierarchical Clustering, lalu mengambil satu perwakilan dari setiap cluster.

\subsection{Feature Importance (MDI)}
Menggunakan Random Forest untuk mengukur *Mean Decrease Impurity*. Fitur yang paling banyak mengurangi entropi adalah fitur terbaik.

\part{Metodologi dan Algoritma}

\chapter{Regime Labeling}
Langkah paling kritis dalam supervised learning bukanlah modelnya, tapi labelnya ("Ground Truth").

\section{Metode Labeling Manual vs Data-Driven}
\subsection{Price-based Rules (Naive)}
"Jika harga di atas SMA200, label = Bull". Terlalu sederhana dan lagging.

\subsection{Drawdown-based Labeling}
Mendifinisikan rezim berdasarkan "rasa sakit" (Drawdown). Berguna untuk analisis risiko historis.

\section{Statistical Change Point Detection}
Mencari titik di mana properti statistik data berubah secara tiba-tiba menggunakan uji statistik (seperti CUSUM atau Chow Test).

\section{Hidden Markov Models (HMM)}
Metode emas (\textit{de facto standard}) untuk pemodelan rezim tanpa pengawasan (\textit{unsupervised}).

\subsection{Konsep Latent States}
HMM berasumsi bahwa pasar digerakkan oleh "mesin keadaan" (\textit{state machine}) yang tidak terlihat. Kita hanya bisa melihat Observasi ($O_t$: return, volatilitas), tapi tidak bisa melihat State ($S_t$: Bull, Bear).

\subsection{Matriks Transisi}
Probabilitas perpindahan antar state. Sifat "lengket" (\textit{sticky}) dari regim pasar dimodelkan di sini (peluang tetap di state yang sama tinggi).

\subsection{Gaussian Mixture Emissions}
Setiap state memancarkan data berdasarkan distribusi probabilitas yang unik.
\begin{itemize}
    \item \textbf{Bull}: Mean Positif, Varians Rendah.
    \item \textbf{Bear}: Mean Negatif, Varians Tinggi.
\end{itemize}

\section{Regime-Switching Models Lainnya}
Selain HMM, kita bisa menggunakan Threshold Auto-Regressive (TAR) atau Markov Switching Auto-Regressive (MSAR).

\chapter{Classification Models untuk Regime Prediction}
Once we have our labels (from HMM or Manual), we can train supervised classifiers to predict them out-of-sample.

\section{Traditional Machine Learning}
\subsection{Logistic Regression}
Model dasar yang memberikan probabilitas linear.

\subsection{Random Forest}
Ensemble dari Decision Trees. Tahan terhadap overfitting dan outlier. Bisa menangani hubungan non-linear antar fitur.

\subsection{Gradient Boosting (XGBoost, LightGBM)}
Model paling populer di kompetisi Kaggle. Sangat akurat tetapi perlu tuning hyperparameter yang hati-hati.

\subsection{Support Vector Machines (SVM)}
Mencari hyperplane pemisah optimal. Bagus untuk dataset kecil dengan dimensi tinggi.

\section{Handling Imbalanced Data}
Data "Crash" sangat jarang.
\begin{itemize}
    \item \textbf{SMOTE}: Oversampling kelas minoritas secara sintetis.
    \item \textbf{Class Weights}: Memberi penalti lebih besar jika salah menebak kelas minoritas.
\end{itemize}

\chapter{Deep Learning untuk Regime Detection}
\section{Feedforward Neural Networks (MLP)}
Jaringan saraf standar. Cocok untuk data tabular. Membutuhkan normalisasi input (Z-score) dan teknik regularisasi (Dropout).

\section{Convolutional Neural Networks (CNN)}
Biasanya untuk gambar, tapi efektif untuk time series 1D. Filter konvolusi bisa belajar mengenali pola lokal seperti "Head and Shoulders" atau lonjakan volatilitas pendek tanpa kita definisikan rumusnya.

\subsection{2D CNN dan Multi-Scale Fusion}
Pendekatan mutakhir (seperti dalam MDPI 2025) menggunakan \textbf{Multi-Scale Fusion} untuk mendeteksi rezim. Caranya adalah dengan mengonversi data return menjadi citra 2D menggunakan \textbf{Gramian Angular Field (GAF)} pada berbagai skala waktu (misalnya 20, 60, dan 120 hari). 

Masing-masing skala ini dimasukkan ke cabang CNN paralel untuk mengekstrak fitur tekstur yang berbeda (mikro vs makro). Fitur-fitur ini kemudian digabungkan (\textit{fusion layer}) untuk menentukan label rezim akhir. Keunggulan metode ini adalah kemampuannya menangkap "konteks hirarkis", di mana model bisa memahami jika volatilitas tinggi saat ini hanyalah kebisingan jangka pendek atau bagian dari transisi rezim jangka panjang.

\section{Recurrent Neural Networks (RNN)}
\subsection{LSTM (Long Short-Term Memory)}
Dirancang untuk mengingat dependensi jangka panjang. LSTM memiliki "memory cell" yang bisa menyimpan informasi tren dari masa lalu.

\subsection{Attention Mechanisms dan Transformers}
\textit{Self-Attention} memungkinkan model untuk "fokus" pada titik-titik waktu tertentu di masa lalu yang paling relevan. Arsitektur \textit{Transformer} (seperti yang digunakan di LLM) kini mulai mendominasi time-series forecasting (misal: Temporal Fusion Transformer).

\chapter{Time Series Models}
\section{Classical Time Series}
\subsection{ARIMA}
Model parametrik klasik. Bagus untuk data stasioner tetapi gagal menangkap perubahan rezim yang tiba-tiba.

\subsection{GARCH untuk Volatility}
Standar industri untuk memprediksi volatilitas. GARCH memodelkan \textit{volatility clustering}.

\subsection{VAR (Vector Autoregression)}
Perluasan multivariate dari ARIMA. Melihat hubungan timbal balik antar beberapa aset.

\section{Modern Time Series ML}
\subsection{Prophet}
Library dari Facebook. Kuat menangani musiman dan liburan.
\subsection{N-BEATS}
Arsitektur Deep Learning murni untuk forecasting yang mengalahkan model statistika.

\chapter{Unsupervised Learning untuk Regime Discovery}
\section{Clustering Methods}
\subsection{K-Means}
Mempartisi data menjadi K kelompok berdasarkan kemiripan fitur (volatilitas, return).
\subsection{GMM (Gaussian Mixture Models)}
Versi probabilistik dari K-Means. GMM mengasumsikan data dihasilkan dari campuran beberapa distribusi Gaussian (Normal). Ini sangat cocok untuk pasar karena kita bisa memodelkan rezim Bull sebagai Gaussian dengan mean positif + varians rendah, dan Bear sebagai Gaussian dengan mean negatif + varians tinggi.

\subsection{Dari GMM ke HMM: Menambahkan Logika Temporal}
Kelemahan utama GMM adalah ia mengabaikan urutan waktu (\textit{sequence}). GMM memperlakukan setiap hari perdagangan sebagai titik data independen. Dalam realitas pasar, jika hari ini adalah Bull Market, kemungkinan besar besok juga tetap Bull Market (sifat \textit{persistence}).
\textbf{Hidden Markov Models (HMM)} menyempurnakan GMM dengan menambahkan \textbf{Matriks Transisi}. HMM tidak hanya melihat "di cluster mana titik ini berada" tetapi juga "seberapa mungkin kita berpindah dari cluster A ke cluster B". Inilah yang membuat HMM menjadi standar emas untuk deteksi rezim yang stabil dan tidak bergejolak (\textit{less noisy signals}).

\section{Dimensionality Reduction}
\subsection{PCA}
Komponen utama pertama seringkali merepresentasikan Beta Pasar.
\subsection{Autoencoders}
Jaringan saraf untuk kompresi data. \textit{Reconstruction Error} dari Autoencoder bisa digunakan sebagai sinyal anomali pasar (Crash Detector).

\chapter{Reinforcement Learning}
\section{Konsep RL untuk Trading}
RL adalah tentang belajar mengambil tindakan (\textit{Action}) untuk memaksimalkan keuntungan jangka panjang (\textit{Reward}).

\section{RL Algorithms}
\subsection{Deep Q-Networks (DQN)}
Menggunakan Neural Network untuk memperkirakan nilai setiap tindakan.
\subsection{PPO (Proximal Policy Optimization)}
Algoritma policy-gradient yang stabil dan menjadi standar industri saat ini untuk trading bot.

\section{RL untuk Regime-aware Trading}
Memasukkan "Label Rezim" (dari HMM) ke dalam state space agen RL, memberinya konteks makro untuk mengambil keputusan yang lebih bijak (misal: lebih konservatif saat Rezim Bear).

\part{Evaluasi dan Optimisasi}

\chapter{Evaluasi Model}
\section{Classification Metrics}
\subsection{Precision, Recall, F1}
Akurasi menyesatkan di finance. Recall (kemampuan mendeteksi Crash) seringkali lebih penting daripada Precision.
\subsection{Confusion Matrix}
Menganalisis tipe kesalahan: False Positive (Alarm Palsu) vs False Negative (Gagal Deteksi Bahaya).

\section{Financial Metrics}
\subsection{Sharpe Ratio & Sortino Ratio}
Mengukur return per unit risiko.
\subsection{Maximum Drawdown}
Resiko kerugian terbesar dari puncak ke lembah.

\section{Backtesting yang Benar}
\subsection{Walk-Forward Validation}
Melatih pada masa lalu, menguji pada masa depan, lalu menggeser jendela waktu.
\subsection{Purged Cross-Validation}
Menghapus data buffer di perbatasan train/test untuk mencegah kebocoran informasi.

\chapter{Model Optimization}
\section{Hyperparameter Tuning}
\subsection{Grid Search vs Random Search}
Random search seringkali lebih efisien.
\subsection{Bayesian Optimization (Optuna)}
Menggunakan probabilitas untuk menebak kombinasi parameter terbaik berikutnya, secara drastis mengurangi waktu pencarian.

\section{Feature Selection}
\subsection{Recursive Feature Elimination (RFE)}
Secara iteratif membuang fitur terlemah.
\subsection{SHAP Values}
Menggunakan Game Theory untuk menjelaskan kontribusi setiap fitur terhadap prediksi model (Explainable AI).

\chapter{Model Interpretability}
\section{Explainable AI (XAI)}
Regulasi menuntut penjelasan. Mengapa model menjual? 
\subsection{SHAP & LIME}
Visualisasi kontribusi fitur lokal dan global.
\subsection{Feature Importance Plot}
Penting untuk sanity check: Apakah fitur yang dipakai model masuk akal secara ekonomi?

\section{Visualisasi Strategis: Regime Overlays}
Salah satu cara terbaik untuk menginterpretasikan model rezim bagi trader adalah melalui \textit{Price Chart Overlays}. Alih-alih hanya melihat angka probabilitas, kita mewarnai latar belakang grafik harga aset (misal: Hijau untuk Bull, Merah untuk Bear, Abu-abu untuk Sideways). Visualisasi ini membantu kita memvalidasi secara kualitatif: "Apakah model berhasil menangkap krisis 2008?" atau "Apakah model terjebak \textit{whipsaw} saat pasar mendatar?".

\chapter{Implementasi Sistem Real-time}
\section{Arsitektur Sistem}
\subsection{Data Pipeline}
Sistem harus memiliki ETL (\textit{Extract, Transform, Load}) otomatis yang mengambil data baru, membersihkannya, dan menghitung fitur sama persis seperti saat fase training (\textit{pre-processing consistency}).

\subsection{Model Serving}
Cara model diakses. Bisa berupa \textit{REST API} (Flask/FastAPI) yang menerima JSON data pasar dan mengembalikan prediksi, atau sistem \textit{embedded} di dalam platform trading.

\subsection{API Design}
Desain endpoint yang efisien (misal: `/predict_regime`, `/get_metrics`). Keamanan (API Key, rate limiting) sangat krusial.

\subsection{Scalability}
Sistem harus mampu menangani lonjakan beban data saat volatilitas pasar tinggi tanpa \textit{latency} yang meningkat.

\section{Real-time Data Processing}
\subsection{Streaming Data}
Menggunakan Kafka atau RabbitMQ untuk menangani aliran data tick-by-tick secara asinkron.

\subsection{Low-latency Processing}
Optimisasi kode Python (misal: menggunakan Cython atau Numba) atau menulis ulang bagian kritis dalam C++/Rust/Go untuk eksekusi order milidetik.

\subsection{Data Quality Monitoring}
Peringatan otomatis jika data feed mati atau memberikan nilai yang tidak masuk akal (misal: harga saham tiba-tiba nol).

\section{Model Deployment}
\subsection{Model Serialization}
Menyimpan model terlatih sebagai file (Pickle, Joblib, ONNX) agar bisa dimuat ulang di server produksi tanpa retraining.

\subsection{Version Control}
Mencatat versi model (misal: v1.0, v1.1) dan data yang digunakan untuk melatihnya (MLflow, DVC). Ini penting untuk audit jika performa tiba-tiba anjlok.

\subsection{A/B Testing}
Menjalankan model baru (Challenger) diam-diam di samping model lama (Champion) dengan uang kertas (\textit{paper trading}) sebelum menggantinya secara live.

\subsection{Model Monitoring}
Dashboard (Grafana) untuk memantau "kesehatan" model: apakah distribusi prediksi berubah? Apakah akurasi harian menurun drastis?

\section{Infrastructure}
\subsection{Cloud vs On-premise}
AWS/GCP menawarkan fleksibilitas dan layanan ML terkelola (\textit{SageMaker}), tapi server on-premise (colocation) lebih cepat (low latency) dan lebih aman untuk strategi HFT.

\subsection{Containerization (Docker)}
Membungkus kode dan semua dependensi dalam kontainer agar berjalan sama persis di laptop developer dan server produksi.

\subsection{Orchestration (Kubernetes)}
Mengelola ratusan kontainer secara otomatis, memastikan sistem selalu hidup (\textit{high availability}).

\chapter{Trading Strategies berbasis Regime}
\section{Regime-conditional Strategies}
Ide inti: Ubah "mode" trading berdasarkan prediksi rezim.

\subsection{Bull Market Strategies}
Fokus: \textbf{Long-only}, \textbf{Momentum}, \textbf{Leverage}.
\begin{itemize}
    \item Beli saat breakout atau pullback (Buy the Dip).
    \item Gunakan leverage moderat karena tren naik cenderung stabil.
    \item Pegang posisi lebih lama (\textit{Let profits run}).
\end{itemize}

\subsection{Bear Market Strategies}
Fokus: \textbf{Cash Conservation}, \textbf{Short Selling}, \textbf{Volatility}.
\begin{itemize}
    \item Keluar dari pasar (100\% Cash) adalah posisi yang valid.
    \item Short-selling indeks atau saham lemah.
    \item Beli volatilitas (Long VIX).
    \item Strategi Mean-Reversion sangat agresif (jual saat rally).
\end{itemize}

\subsection{Sideways Market Strategies}
Fokus: \textbf{Income Generation}, \textbf{Mean Reversion}.
\begin{itemize}
    \item Jual opsi (Covered Call, Iron Condor) untuk mendapatkan premium dari time decay (\textit{Theta}).
    \item Beli di Support bawah, Jual di Resistance atas (\textit{Range Trading}).
\end{itemize}

\section{Portfolio Allocation}
\subsection{Dynamic Asset Allocation}
Menggeser bobot portfolio dari Saham (Equity) ke Obligasi/Emas saat probabilitas rezim resesi > 50\%.

\subsection{Risk Parity}
Mengalokasikan aset berdasarkan kontribusi risiko, bukan nilai dolar. Di rezim volatilitas tinggi, posisi saham dikurangi secara otomatis untuk menjaga target volatilitas portfolio konstan.

\subsection{Tactical Allocation}
Perubahan jangka pendek untuk memanfaatkan inefisiensi sesaat (misal: rotasi sektor dari Tech ke Energy saat sinyal inflasi menyala).

\section{Risk Management}
\subsection{Position Sizing}
Rumus Kelly Criterion yang disesuaikan dengan keyakinan rezim. Pertaruhkan lebih sedikit saat sinyal rezim tidak pasti.

\subsection{Stop Loss Strategies}
Stop loss longgar di Bull Market (biarkan bernafas), Stop loss sangat ketat di Bear Market.

\subsection{Hedging}
Membeli asuransi portofolio (Options Put) hanya saat sinyal "Bahaya" menyala, sehingga menghemat biaya premi asuransi saat pasar aman.

\subsection{Diversification}
Di saat krisis, korelasi aset mendekati 1. Diversifikasi tradisional gagal. Strategi rezim mencari aset "Crisis Alpha" (seperti Managed Futures) yang justru naik saat saham jatuh.

\section{Execution}
\subsection{Order Types}
Market order untuk masuk cepat saat breakout rezim penting. Limit order untuk mengambil profit di target.

\subsection{Execution Algorithms}
VWAP/TWAP algo untuk memecah order besar agar tidak menggerakkan pasar.

\subsection{Market vs Limit Orders}
Kapan harus agresif (memakan likuiditas) dan kapan harus pasif (menyediakan likuiditas dan dapat rebate).

\chapter{Studi Kasus: Saham}
\section{Data Preparation}
Menggunakan data harian S\&P 500 (SPY) dari tahun 2000-2023. Fitur mencakup technical indicators (RSI, MACD) dan data makro (Yield Curve 10Y-2Y, VIX). Data dibagi menjadi Train (2000-2015), Validation (2016-2019), dan Test (2020-2023).

\section{Feature Engineering}
Membuat fitur "Jarak dari All-Time High" dan "Sektor Rotasi Alpha" (XLK vs XLE). Label rezim dibuat menggunakan algoritma HMM 2-state pada volatilitas mingguan.

\section{Model Development}
Melatih Random Forest Classifier dengan `class_weight='balanced'`. Hyperparameter `max_depth` dibatasi 5 untuk mencegah overfitting.

\section{Backtesting Results}
Strategi "Long saat Bull, Cash saat Bear" menghasilkan CAGR 12\% vs 8\% Buy-and-Hold, dengan Max Drawdown hanya -15\% (vs -34% SPY). Rasio Sharpe meningkat dari 0.6 menjadi 1.1.

\section{Live Trading Performance}
Paper trading selama 6 bulan menunjukkan slippage 0.05\% per trade. Model berhasil keluar sebelum koreksi pasar Agustus 2023.

\section{Lessons Learned}
Rezim "Sideways" seringkali salah diklasifikasikan sebagai awal Bull, menyebabkan "whipsaw" (beli lalu segera rugi). Penambahan filter ADX membantu mengurangi signal palsu ini.

\chapter{Studi Kasus: Cryptocurrency}
\section{Karakteristik Crypto Markets}
Pasar 24/7 tanpa henti, volatilitas ekstrim (sering >5\% sehari), dan dominasi emosi ritel. Data on-chain (aktivitas dompet paus) adalah sumber alpha unik.

\section{Data Sources}
Data OHLCV 1-jam dari Binance untuk BTCUSDT dan ETHUSDT. Data sentimen dari "Fear and Greed Index" Crypto. Data on-chain dari Glassnode (Net Unrealized Profit/Loss).

\section{Model Adaptation}
Menggunakan LSTM (Deep Learning) karena pola sekuensial "pump and dump" sangat kuat. Window input diperpendek menjadi 60 jam (vs 60 hari di saham) karena dinamika pasar yang lebih cepat.

\section{Results dan Analysis}
Model 3-state (Bull, Bear, Crash). Strategi berhasil menghindari crash FTX (Nov 2022) dengan beralih ke stablecoin 2 hari sebelumnya karena mendeteksi lonjakan volatilitas aneh dan aliran keluar dana besar-besaran. Namun, model sering terlambat masuk kembali saat "V-shape recovery".

\chapter{Studi Kasus: Forex}
\section{Currency Pair Selection}
Fokus pada pasangan mata uang utama (EURUSD, USDJPY, GBPUSD) yang memiliki likuiditas tertinggi dan spread terendah.

\section{Multi-timeframe Analysis}
Model menggunakan input dari grafik H4 (4 jam) untuk sinyal taktis, dan grafik D1 (Harian) untuk bias strategis.

\section{Macro Factors}
Memasukkan selisih suku bunga (Interest Rate Differential) antara Fed dan ECB sebagai fitur utama. 

\section{Implementation}
Strategi "Mean Reversion" aktif saat sesi Asia (sideways), dan strategi "Momentum Breakout" aktif saat sesi London/New York overlap. Model ML bertindak sebagai "dosen" yang memilih strategi mana yang dipakai jam ini.

\chapter{Studi Kasus: Multi-asset}
\section{Cross-asset Correlations}
Menganalisis matriks korelasi bergulir antara Saham, Obligasi, Emas, dan Bitcoin.

\section{Global Regime Detection}
Membangun "Global Risk Indicator". Jika semua aset berkorelasi positif dan jatuh bersamaan, ini adalah rezim "Liquidity Crisis" (seperti Maret 2020).

\section{Portfolio Construction}
Menggunakan "Hierarchical Risk Parity" (HRP) yang ditenagai oleh prediksi rezim untuk mengalokasikan bobot aset.

\section{Performance Evaluation}
Portfolio multi-aset ini memberikan kurva ekuitas paling mulus, dengan volatilitas tahunan di bawah 8\% meskipun return setara saham murni.

\part{Topics Lanjutan}

\chapter{Online Learning dan Adaptive Models}
\section{Concept Drift}
\subsection{Deteksi Drift}
Pasar berevolusi. Model yang dilatih pada data 2010 mungkin gagal total di 2024. Kita menggunakan uji statistik (seperti Kolmogorov-Smirnov test) untuk mendeteksi apakah distribusi data masuk baru berbeda signifikan dari data latih.

\subsection{Adaptation Strategies}
Jika drift terdeteksi, kita punya dua pilihan: Melatih ulang model (\textit{retraining}) atau memberikan bobot lebih pada data terbaru (\textit{instance weighting}).

\section{Online Learning Algorithms}
\subsection{Incremental Learning}
Algoritma yang bisa belajar satu per satu data baru tanpa melupakan ilmu lama sepenuhnya, dan tanpa perlu melatih ulang dari nol (misalnya: Online Random Forest).

\subsection{Transfer Learning}
Melatih model dasar pada data pasar AS yang melimpah, lalu "mentransfer" pengetahuan itu untuk memprediksi pasar berkembang (seperti Indonesia) yang datanya sedikit.

\subsection{Meta-learning}
Belajar untuk beradaptasi dengan cepat. Model dilatih pada berbagai tugas (saham, crypto, forex) sehingga ketika menghadapi aset baru, ia bisa belajar hanya dengan sedikit contoh (\textit{few-shot learning}).

\section{Model Updating}
\subsection{Retraining Frequency}
Seberapa sering kita melatih ulang? Setiap malam? Setiap minggu? Trade-off antara biaya komputasi dan kesegaran model.

\subsection{Sliding Window}
Menggunakan jendela data geser (misal: 1000 hari terakhir). Data yang lebih tua dari itu dibuang karena dianggap "basi".

\subsection{Weighted Recent Data}
Memberikan bobot eksponensial pada sampel data. Data kemarin jauh lebih penting daripada data 5 tahun lalu.

\chapter{Alternative Approaches}
\section{Bayesian Methods}
\subsection{Bayesian Networks}
Grafik asiklik berarah yang merepresentasikan hubungan probabilitas bersyarat. "Jika Minyak naik, maka inflasi naik (80\%), maka Suku Bunga naik (70\%), maka Saham turun."

\subsection{Gaussian Processes}
Metode non-parametrik yang, alih-alih memberikan satu garis prediksi, memberikan "awan probabilitas". Sangat bagus untuk mengukur ketidakpastian.

\subsection{Bayesian Optimization}
Efisien untuk mencari parameter strategi trading (misal: panjang MA) dengan jumlah percobaan minimal.

\section{Fuzzy Logic}
\subsection{Fuzzy Sets}
Pasar tidak hitam putih. Daripada "Bull" atau "Bear", logika fuzzy memungkinkan "Agak Bullish" atau "Sangat Bearish".

\subsection{Fuzzy Rules}
"JIKA harga naik CEPAT DAN volume TINGGI, MAKA BELI KUAT". Aturan linguistik ini lebih mudah dipahami manusia daripada bobot Neural Net.

\subsection{Fuzzy Inference Systems}
Sistem pengambilan keputusan berbasis aturan-aturan fuzzy.

\section{Genetic Algorithms}
\subsection{Evolution Strategies}
Membiakkan ribuan strategi trading acak. Yang rugi "dibunuh", yang untung "dikawinkan" (crossover parameternya) untuk menghasilkan generasi anak yang lebih baik.

\subsection{Strategy Optimization}
Menemukan kombinasi parameter indikator yang optimal tanpa terjebak di optimum lokal.

\section{Graph Neural Networks (GNN)}
\subsection{Market Networks}
Memodelkan pasar sebagai grafik raksasa di mana simpul adalah Saham dan sisi adalah korelasi. GNN bisa memprediksi harga saham A berdasarkan pergerakan tetangga-tetangganya di grafik.

\subsection{Correlation Graphs}
Struktur grafik bisa berubah dinamis setiap hari, menggambarkan dinamika penularan risiko sistemik.

\chapter{Explainability dan Transparency}
\section{Black Box vs White Box Models}
Regulator keuangan membenci Black Box. Kita harus bisa menjawab: "Mengapa algoritma menolak kredit nasabah ini?" atau "Mengapa bot menjual semua saham hari ini?"

\section{Regulatory Requirements}
GDPR di Eropa menuntut "Right to Explanation". Model deep learning murni seringkali gagal memenuhi standar ini.

\section{Stakeholder Communication}
Manajer investasi harus bisa menjelaskan logika model kepada investor. "Algoritma bilang begitu" bukan jawaban yang bisa diterima saat rugi.

\section{Trust dan Validation}
Kepercayaan dibangun melalui ribuan trade yang konsisten dan penjelasan yang masuk akal atas setiap keputusan.

\chapter{Trends dan Inovasi}
\section{Quantum Machine Learning}
Komputer kuantum menjanjikan percepatan eksponensial dalam optimisasi portofolio dan simulasi Monte Carlo.

\section{Federated Learning}
Bank-bank bisa melatih model deteksi penipuan bersama-sama tanpa perlu saling berbagi data nasabah sensitif. Privasi terjaga, tapi model jadi lebih pintar.

\section{Neuromorphic Computing}
Chip komputer yang meniru struktur otak manusia, memungkinkan AI berjalan dengan daya sangat rendah di perangkat edge (misal: satelit).

\section{Alternative Data Explosion}
Data semakin aneh: memantau bayangan tangki minyak dari luar angkasa, melacak kapal kargo via AIS, hingga menganalisis nada suara CEO saat earnings call.

\chapter{Challenges dan Limitations}
\section{Market Regime Changes}
Tantangan abadi. Segera setelah kita menemukan pola untung, pasar beradaptasi dan pola itu hilang (\textit{Alpha Decay}).

\section{Black Swan Events}
ML dilatih pada data historis. Ia tidak bisa memprediksi kejadian yang belum pernah terjadi sebelumnya (seperti Pandemi COVID-19 pertama kali).

\section{Data Quality Issues}
"Garbage In, Garbage Out". Sebagus apapun algoritmanya, jika data tick-nya bolong-bolong, hasilnya sampah.

\section{Computational Limits}
Melatih Transformer raksasa membutuhkan biaya listrik dan GPU yang masif, seringkali di luar jangkauan trader ritel.

\section{Ethical Considerations}
Apakah etis menggunakan algoritma untuk memprediksi kebangkrutan negara dan mengambil untung darinya? Atau menggunakan HFT untuk, secara teknis, "mencuri" sen dari investor lambat (front-running)?

\chapter{Kesimpulan}
\section{Rangkuman}
Perjalanan kita dari teori pasar efisien, menembus kekacauan data tick, hingga arsitektur Deep Learning canggih.

\section{Best Practices}
Validasi yang ketat (\textit{Purged Cross-Validation}), manajemen risiko paranoik, dan kerendahan hati bahwa pasar selalu benar.

\section{Practical Recommendations}
Mulailah dari yang sederhana (Logistic Regression). Fokus 80\% waktu pada data, bukan model. Jangan pernah percaya hasil backtest 100\%.

\section{Final Thoughts}
Machine Learning di finance bukan tentang menggantikan manusia, tapi memperkuat intuisi manusia dengan kemampuan komputasi statistik skala besar. Trader masa depan adalah Centaur: setengah manusia, setengah AI.

\appendix

\chapter{Kode dan Implementasi}
\section{Setup Environment}
\subsection{Python Libraries}
Pandas, NumPy, Scikit-learn, PyTorch, TensorFlow, TA-Lib, Zipline/Backtrader.

\subsection{R Packages}
quantmod, PerformanceAnalytics, TTR.

\subsection{Trading Platforms}
MetaTrader 5 (MQL5), Interactive Brokers API, CCXT (untuk Crypto).

\section{Data Acquisition}
\subsection{API Examples}
Kode contoh menarik data dari Yahoo Finance (`yfinance`) dan Binance API.

\subsection{Web Scraping Code}
Skrip BeautifulSoup sederhana untuk mengambil data sentimen, data alternatif.

\subsection{Database Setup}
Menyimpan data time series di PostgreSQL dengan TimescaleDB atau InfluxDB.

\section{Model Templates}
\subsection{Classification Template}
Scikit-learn pipeline standar (Scaler -> PCA -> Random Forest).

\subsection{Deep Learning Template}
Kelas PyTorch Lightning untuk LSTM.

\subsection{Backtesting Template}
Kerangka kerja `Backtrader` untuk menjalankan strategi regime-switching.

\subsection{HMM Implementation Template (hmmlearn)}
Contoh kode menggunakan library \texttt{hmmlearn} untuk deteksi rezim:
\begin{verbatim}
from hmmlearn.hmm import GaussianHMM
import yfinance as yf

# Preprocessing
data = yf.download("SPY", start="2010-01-01")
returns = np.log(data['Close'] / data['Close'].shift(1)).dropna()

# Fit Model
model = GaussianHMM(n_components=2, covariance_type="full", n_iter=100)
model.fit(returns.values.reshape(-1, 1))

# Predict Regimes
hidden_states = model.predict(returns.values.reshape(-1, 1))
\end{verbatim}

\section{Complete Examples}
\subsection{End-to-end Pipeline}
Notebook Jupyter lengkap dari download data sampai evaluasi Sharpe Ratio.

\subsection{Production-ready Code}
Struktur folder proyek Python yang bersih dengan logging dan error handling.

\chapter{Glossary}
\section{Machine Learning Terms}
Bias, Variance, Overfitting, Epoch, Gradient Descent, Hyperparameter.

\section{Finance Terms}
Alpha, Beta, Delta, Gamma, Theta, Sharpe Ratio, Drawdown, Slippage.

\section{Statistical Terms}
Stationarity, Autocorrelation, Heteroskedasticity, Kurtosis, Skewness.

\backmatter

\chapter*{Daftar Pustaka}
\addcontentsline{toc}{chapter}{Daftar Pustaka}
[1] Lopez de Prado, M. (2018). \textit{Advances in Financial Machine Learning}. Wiley. \\
[2] Aronson, D. R. (2007). \textit{Evidence-Based Technical Analysis}. Wiley. \\
[3] Jansen, S. (2020). \textit{Machine Learning for Algorithmic Trading}. Packt.

\chapter*{Indeks}
\addcontentsline{toc}{chapter}{Indeks}

\vfill
\begin{center}
    \textit{Terakhir diperbarui: 4 Februari 2026, 16:58 WIB}
\end{center}

\end{document}
